\section{Profilage et validation des performances}

Afin d'isoler avec précision les goulots d'étranglement du système et de valider l'architecture proposée, une 
série de mesures de performances a été menée. L'objectif est de quantifier le temps de transit d'une image 
depuis le capteur jusqu'à l'antenne radio (Latence TX), puis son temps de traitement à la réception 
(Latence RX), afin d'établir un bilan "Glass-to-Glass".

\subsection{Profilage de la chaîne d'émission (TX)}

\subsubsection{Temps d'encodage matériel (Jetson)}
Pour mesurer le temps de traitement de l'encodeur matériel Nvidia (\texttt{nvv4l2h264enc}), les outils de 
traçage natifs de GStreamer\index{GStreamer} ont été employés. L'activation de la variable d'environnement 
\texttt{GST\_TRACERS="latency(flags=element)"} permet d'insérer des sondes (probes) automatiques entre chaque 
élément du pipeline. 

L'analyse des journaux d'exécution a démontré que le composant matériel traite et compresse une trame vidéo 
en un temps extrêmement stable d'environ \textbf{1.2 ms}.

\subsubsection{Mesure au cycle d'horloge du traitement MCU (STM32)}
Pour profiler le temps mis par le microcontrôleur pour transférer un bloc de 1400 octets depuis le buffer 
SPI\index{SPI} vers la pile réseau LwIP\index{LwIP}, les fonctions classiques de type \texttt{HAL\_GetTick()} 
(résolution de l'ordre de la milliseconde) se sont avérées insuffisantes. 

Le code exploite donc le composant matériel ARM CoreSight DWT\index{DWT} (\textit{Data Watchpoint and Trace}). 
Le registre \texttt{DWT->CYCCNT} s'incrémente à chaque coup d'horloge du processeur. Le STM32 fonctionnant à 
160 MHz, un cycle correspond à une résolution absolue de 6.25 nanosecondes.

\begin{lstlisting}[language=C]
// Demarrage du chronometre DWT dans l'interruption SPI
dwt_start = DWT->CYCCNT;

// ... [Encapsulation pbuf_alloc et envoi udp_sendto] ...

// Arret du chronometre et calcul du temps en microsecondes
dwt_end = DWT->CYCCNT;
mcu_cycles = dwt_end - dwt_start;
uint32_t freq_mhz = SystemCoreClock / 1000000; // 160 MHz
uint32_t time_us = mcu_cycles / freq_mhz;
\end{lstlisting}

L'extraction de cette métrique a révélé un temps de traitement MCU ultra-court de \textbf{0.8 ms par paquet} 
(soit environ 128 000 cycles d'horloge). Le STM32 agit donc comme un pont réseau "zéro-délai".

\subsubsection{Bilan de la latence d'émission par image}
Le temps d'émission d'une image complète est sérialisé par le \textit{Handshake} matériel. À un bitrate de 
400 kbps et 24 FPS, une image pèse environ 2083 octets, nécessitant la transmission de 2 paquets de 1400 octets. 
Le transfert SPI théorique à 5 MHz d'un paquet de 1400 octets prend 2.24 ms.

Le temps total d'émission d'une image compressée est donc de :
$$Latence_{TX} = T_{Encodeur} + \left( Nb_{Paquets} \times (T_{SPI} + T_{MCU}) \right)$$
$$Latence_{TX} = 1.2\text{ ms} + \left( 2 \times (2.24\text{ ms} + 0.8\text{ ms}) \right) = \mathbf{7.28\text{ ms}}$$
Ce temps d'à peine 7 millisecondes confirme l'excellence de l'architecture matérielle embarquée.

\subsection{Profilage de la réception (RX)}
Du côté de la station au sol embarquée (Raspberry Pi 4), le recours au décodage logiciel (\texttt{avdec\_h264}) 
a été privilégié pour contourner la file d'attente du décodeur matériel.
L'utilisation des \texttt{GST\_TRACERS} sur la station de réception a prouvé que le processeur ARM de la 
Raspberry Pi parvient à décoder une trame entrante en une moyenne de \textbf{3.0 ms}, avec un affichage 
quasi-instantané assuré par l'élément bas niveau \texttt{kmssink}.

\subsection{Bilan final "Glass-to-Glass"}
La latence "Glass-to-Glass" absolue (le délai entre un événement réel et son affichage sur l'écran du pilote) 
a été mesurée de manière empirique en filmant un chronomètre à haute fréquence de rafraîchissement.

La valeur stabilisée de cette latence se situe entre \textbf{160 ms et 200 ms}.

Ce délai s'explique par la somme des traitements matériels incompressibles, principalement dominés par le temps 
d'acquisition (la durée physique d'une trame), le buffer de l'ISP (Image Signal Processor) de la caméra 
(évalué à environ 133 ms pour un échantillonnage à 30 FPS), et les temps de propagation radio de la modulation 
Sub-1 GHz.

Bien que cette latence soit supérieure aux 20-30 ms des systèmes FPV analogiques purs (5.8 GHz), elle reste 
suffisamment basse et stable pour assurer la navigation d'un drone, tout en apportant l'avantage décisif d'une 
liaison radio robuste capable de traverser d'importants obstacles physiques (NLOS).